%% ---------------------------------------------------------------------
%% 8. GETTING STARTED — Required · 1–2 pages
%% The first 2–4 weeks: how the student gets acquainted with the project.
%% Include:
%%  - first-meeting checklist (accounts, repository, lab access, equipment handover);
%%  - environment setup steps (tool versions, install guides, a "hello world");
%%  - 2–4 WARM-UP EXERCISES that build toward the project, each with a clear
%%    "done when" acceptance condition and an estimated time:
%%      \begin{warmup}{Title}{6 h} ... \donewhen{...} \end{warmup}
%%  - what to read first (point to the most important rows of ch. 7).
%% The first checkpoint (ch. 9) should be the completion of these exercises.
%% ---------------------------------------------------------------------
\chapter{Getting Started}
\label{ch:start}

\section{First meeting checklist}
\begin{itemize}
  \item \missing{repository / shared folder access}
  \item \missing{lab / equipment access, safety briefing}
  \item \missing{communication channel and meeting slot}
\end{itemize}

\section{Environment setup}
\begin{enumerate}
  \item \missing{install tool X (version), licence instructions}
  \item \missing{build and run the reference example}
\end{enumerate}

\section{Warm-up exercises}
\begin{warmup}{\missing{Exercise title}}{\missing{6 h}}
  \missing{What to do, in 2–4 sentences, with links to the needed resources.}
  \donewhen{\missing{observable acceptance condition, e.g. “simulation shows …”}}
\end{warmup}

\begin{warmup}{\missing{Exercise title}}{\missing{8 h}}
  \missing{Second exercise that builds on the first and is a small piece of the real project.}
  \donewhen{\missing{acceptance condition}}
\end{warmup}

\section{Read first}
\missing{The 2–3 resources to start with, in order.}
